% P16-C3 and a finite identification component of P16-C4.
\chapter{Finite Structural Causal Models}
\label{sec:p16-finite-causal-models}

A structural causal model specifies how its variables are assigned values
and how those assignments change under a specified intervention. Pearl's
nonparametric structural equations and equation-replacement semantics supply
the source presentation used here
\cite[secs.~3.1--3.2.3, pp.~103--112]{Pearl2009CausalOverview}.
This chapter gives a finite formulation with complete proofs, including the
probability assumptions required for factorization. All numerical examples
are stipulated mathematical models.

The probability language is that of
\cref{sec:p15-finite-conditioning,sec:p15-finite-independence}.
The background implications of logical abduction in
\cref{sec:p16-finite-abduction} do not themselves specify the assignments
or interventions defined below.

\section{Graphs, assignments and disturbances}
\label{sec:p16-causal-model-data}

\begin{hbdefinition}{def:p16-causal-dag}{A finite directed acyclic graph}
Let \(N=\{1,\ldots,n\}\), with \(n\geq1\). A \emph{directed graph}
\(G=(N,E)\) consists of the vertices \(N\) and a set
\(E\subseteq N\times N\) of directed edges. Write \(j\to i\)
for \((j,i)\in E\), and put
\[
 \operatorname{pa}_G(i)\defeq\{j\in N:j\to i\}.
\]
These vertices are the \emph{parents} of \(i\); when \(j\to i\), the
vertex \(i\) is a \emph{child} of \(j\). A directed path of
positive length is a sequence \(i_0,\ldots,i_m\), with \(m\geq1\),
such that \(i_{r-1}\to i_r\) for every \(r\). A directed cycle
is such a path with \(i_0=i_m\). The graph is \emph{acyclic} when
there is no directed cycle, including a loop of length one. A finite
directed acyclic graph is abbreviated \emph{DAG}.

For \(I\subseteq N\), let \(\operatorname{de}_G(I)\) be the vertices
reachable from a member of \(I\) by a directed path of positive length.
Its members are the \emph{descendants} of \(I\); a vertex \(j\) is an
\emph{ancestor} of \(i\) when \(i\in\operatorname{de}_G(\{j\})\).
A \emph{topological order} is an enumeration \(i_1,\ldots,i_n\) of
all vertices in which every parent precedes its child.
\index{causal graph!directed acyclic graph}\index{topological order}
\end{hbdefinition}

\begin{hbtheorem}{thm:p16-causal-topological-order}{proved=annex:p16.causal-topological-order}{Topological evaluation is available}
Every finite DAG has a topological order. Repeatedly removing a vertex
with no incoming edge in the remaining graph produces one after exactly
\(n\) removals.
\end{hbtheorem}

\begin{hbdefinition}{def:p16-finite-scm}{A finite structural causal model}
A model \(M\) consists of the following data.
\begin{enumerate}
\item A DAG \(G=(N,E)\), whose vertices index the \emph{endogenous}
      variables to be assigned by the equations, and a finite nonempty
      value set \(S_i\) for each vertex \(i\). For \(J\subseteq N\), put
      \(S_J=\prod_{j\in J}S_j\). Products are indexed tuples;
      \(S_\varnothing\) contains the single empty tuple. If
      \(v\in S_N\), write \(v_J=(v_j)_{j\in J}\).
\item A finite nonempty outcome space \(\Omega\), with a mass
      function \(q:\Omega\to[0,1]\) satisfying
      \(\sum_{\omega\in\Omega}q(\omega)=1\). Its event probability
      is \(\mathbb P(A)=\sum_{\omega\in A}q(\omega)\).
\item For each \(i\), a finite nonempty disturbance value set
      \(\mathcal U_i\), a function \(U_i:\Omega\to\mathcal U_i\),
      and a total structural function
      \[
       f_i:S_{\operatorname{pa}_G(i)}\times\mathcal U_i\longrightarrow S_i.
      \]
\end{enumerate}
The \(U_i\) are \emph{exogenous disturbances}: this model gives their
joint probability law but no equations assigning them from the endogenous
variables. They may share information or be dependent. Their independence
will be an additional assumption when explicitly stated.

At an outcome \(\omega\), a \emph{solution} is an assignment
\(v\in S_N\) satisfying all the structural equations
\[
 v_i=f_i(v_{\operatorname{pa}_G(i)},U_i(\omega)),\qquad i\in N.
\]
The indexed left sides declare which variable each function assigns.
Solving an equation algebraically for another variable does not replace
that declaration. The parent sets state which endogenous inputs a function
is allowed to use. A listed parent may be ignored by the function; an edge
does not assert a nonzero effect or a detectable association.
\index{structural causal model!finite}\index{exogenous disturbance}
\end{hbdefinition}

Pearl introduces these ingredients through Sewall Wright's (1921)
treatment of a disease and its symptom, the first attempt, on his account,
to state mathematically that symptoms do not cause diseases
\cite[sec.~3.1, pp.~103--105]{Pearl2009CausalOverview}. With \(x\) the
level of a disease \(X\) and \(y\) the level of a symptom \(Y\), Wright
writes the linear equation
\[
 y=\beta x+u_Y,
\]
the source's equation~(1), where \(u_Y\) stands for all factors other
than the disease that could affect \(Y\) when \(X\) is held constant; the
disease itself may be ``explained'' by \(x=u_X\), with \(U_X\) standing
for all factors affecting \(X\). The equation is read as an assignment
process: Nature examines the values of \(x\) and \(u\) and gives \(Y\) the
value \(\beta x+u_Y\) (pp.~103--104). The source's note~6 states that the
linear form serves only as an illustration of the historical development
of path analysis, not as a typical disease--symptom relation, and that the
variables are standardized, with mean zero and variance one (p.~103).
Because an algebraic equation is symmetric, (1) can be rewritten as
\[
 x=(y-u_Y)/\beta,
\]
the source's equation~(2), which might be misread as saying that the
symptom influences the disease. Wright's remedy was a diagram, later
called a path diagram, with arrows drawn from perceived causes to their
perceived effects. In the source's Fig.~1 it has the solid arrow
\(X\to Y\) labelled \(\beta\) and dashed arrows \(U_X\to X\) and
\(U_Y\to Y\); in Fig.~1(b) a dashed double arrow joins \(U_X\) and
\(U_Y\), the customary mark that correlation between the two disturbances
is presumed possible, while Fig.~1(a) omits it (pp.~104--105). More
importantly, the source reads an absent arrow as an empirical claim: the
value Nature gives the one variable does not depend on the other. The missing arrow from
\(Y\) to \(X\) says that the symptom is not among the factors \(U_X\)
affecting the disease (p.~104).

The two parts of the model then have distinct roles in the source. The
diagram encodes the possible existence of a direct causal influence of
\(X\) on \(Y\) and the absence of any influence of \(Y\) on \(X\); the
equations encode the quantitative relations among the variables, which
are ``to be determined from the data'' (p.~104). The parameter \(\beta\)
is called a \emph{path coefficient} and quantifies the direct causal
effect of \(X\) on \(Y\). Given \(\beta\) and \(U_Y\), the equation
claims that a unit increase in \(X\) produces an increase of \(\beta\)
units in \(Y\), whatever the values of the other variables in the model,
and the source stresses that this holds whether the increase in \(X\) is
produced from outside the model or from within it (p.~104). In the finite
formulation that last clause is the equation-replacement semantics of
\cref{def:p16-causal-intervention}: an externally set value of \(X\)
enters the unchanged function \(f_Y\) with the unchanged disturbance
\(U_Y\), exactly as an internally produced value does.
\Cref{thm:p15-w02c-p15-009-linear-effect} proves the resulting mean
statements for a linear response equation, including Wright's covariance
relations \(\operatorname{Cov}(X,Y)=\beta\) for Fig.~1(a) and
\(\operatorname{Cov}(X,Y)=\beta+\operatorname{Cov}(U_Y,U_X)\) for
Fig.~1(b), the source's equations~(3) and~(4) (p.~105).

On the disturbances, Pearl's reading is as follows (p.~104). The
variables \(U_X,U_Y\) are \emph{exogenous}: background factors, observed
or not, that the modeler chooses to leave unexplained, which influence the
other variables of the model (the \emph{endogenous} ones) and are not
influenced by them. When unobserved they are also called
``disturbances'' or ``errors'', and they stand for factors left out of
the model but judged relevant to explaining how its variables behave.
They belong to physical reality, for instance genetic factors or
socio-economic conditions, and they are what produces the variation seen
in the data. They are treated like any other variable, but their values
often cannot be measured precisely, so that one must settle for
acknowledging that they exist and judging qualitatively how they relate to
the other variables of the system. They therefore differ fundamentally from
the residual terms of a regression equation, which are artifacts of the
analysis and are uncorrelated with the regressors by definition. Nothing in
\cref{def:p16-finite-scm} forces a disturbance to be uncorrelated with any
endogenous variable; whether disturbances are independent is the separate
assumption stated below.

Path diagrams are read with kinship words, whose meaning the source
takes to be usually self-evident (pp.~104--105): an arrow \(X\to Y\)
makes \(X\) a parent of \(Y\) and \(Y\) a child of \(X\), and the
diagram's ancestors and descendants are understood in the same way. A
\emph{path}, in the source, is any consecutive sequence of edges, solid or
dashed; its example is that Fig.~1(b) has two paths between \(X\) and
\(Y\), the direct arrow \(X\to Y\) and the path that traces the nodes
\(X\), \(U_X\), \(U_Y\) and \(Y\) (p.~105). The following definition
gives these terms for the diagram of a finite model, disturbance vertices
and dashed edges included. Restricted to the endogenous vertices and
solid arrows they are the terms of \cref{def:p16-causal-dag}. The paths
of \cref{def:p16-active-path} are the paths of this definition that use
solid arrows only.

\begin{hbdefinition}{def:p15-w02c-p15-009-causal-diagram}{The causal diagram of a finite model}
Let \(M\) be a model as in \cref{def:p16-finite-scm}, with graph
\(G=(N,E)\). Its \emph{causal diagram} has \(2n\) vertices: the
endogenous vertices \(i\in N\) and, for each \(i\), a \emph{disturbance
vertex}, written \(U_i\) like the disturbance it stands for. Its edges are
of three kinds:
\begin{enumerate}
\item a \emph{solid arrow} \(j\to i\) for each \((j,i)\in E\);
\item a \emph{dashed arrow} \(U_i\to i\) for each \(i\in N\);
\item a \emph{dashed double arrow} \(U_i\leftrightarrow U_j\) for each
      unordered pair \(\{i,j\}\), \(i\ne j\), in a specified set
      \(E^{\leftrightarrow}\), drawn when dependence between \(U_i\) and
      \(U_j\) is allowed.
\end{enumerate}
An arrow \(a\to b\), solid or dashed, makes \(a\) a \emph{parent} of
\(b\) and \(b\) a \emph{child} of \(a\). A directed path follows arrows
in their direction; double arrows have no direction and are not used in
directed paths. The vertex \(b\) is a \emph{descendant} of \(a\), and
\(a\) an \emph{ancestor} of \(b\), when a directed path of positive length
leads from \(a\) to \(b\). For distinct vertices \(a,b\), a \emph{path
between \(a\) and \(b\)} is a sequence of distinct vertices
\(w_0=a,w_1,\ldots,w_m=b\), with \(m\geq1\), in which each consecutive
pair is joined by an edge of one of the three kinds, used in either
direction; the kind and direction of each edge used are retained.

A diagram with \(E^{\leftrightarrow}=\varnothing\) is used in this
chapter only together with mutual independence of the whole disturbance
family (\cref{def:p16-causal-kernels}). When \(E^{\leftrightarrow}\) is
nonempty the model of \cref{def:p16-finite-scm} places no restriction on
the joint law of the disturbances.
\index{causal diagram}\index{path!causal diagram}
\end{hbdefinition}

The requirement that a path's vertices be distinct makes the source's
count exact: if repetitions were allowed, the edge
\(U_X\leftrightarrow U_Y\) could be traversed back and forth, and
Fig.~1(b) would have infinitely many such sequences between \(X\) and
\(Y\). The distinctness requirement is this handbook's precision; the
source does not state it.

\begin{hbtheorem}{thm:p15-w02c-p15-009-diagram-paths}{proved=annex:p15-w02c-p15-009-diagram-paths}{Kinship of disturbance vertices and the paths of Fig.~1}
\begin{enumerate}
\item In the causal diagram of every finite model, a disturbance vertex
      \(U_i\) has no parent, its only child is \(i\), and its descendants
      are the members of \(\{i\}\cup\operatorname{de}_G(\{i\})\). The
      diagram has no directed cycle. For endogenous \(i,j\), the vertex
      \(j\) is a parent, child, ancestor or descendant of \(i\) in the
      diagram exactly when it is one in \(G\) by
      \cref{def:p16-causal-dag}; the disturbance vertices among the
      ancestors of \(i\) are the \(U_k\) with
      \(i\in\{k\}\cup\operatorname{de}_G(\{k\})\).
\item Let \(N=\{X,Y\}\) and \(E=\{(X,Y)\}\). If
      \(E^{\leftrightarrow}=\{\{X,Y\}\}\), the diagram of the source's
      Fig.~1(b), there are exactly two paths between \(X\) and \(Y\):
      \((X,Y)\), using the solid arrow \(X\to Y\), and
      \((X,U_X,U_Y,Y)\), using \(U_X\to X\) against its direction,
      then \(U_X\leftrightarrow U_Y\), then \(U_Y\to Y\). If
      \(E^{\leftrightarrow}=\varnothing\), the diagram of Fig.~1(a), the
      only path between \(X\) and \(Y\) is \((X,Y)\).
\end{enumerate}
\end{hbtheorem}

Wright's major contribution to causal analysis, according to Pearl, was,
besides the language of path diagrams itself, a set of graphical rules
for expressing the covariance between two observed variables through the
path coefficients together with the error-term covariances; the source's
equations~(3) and~(4), \(\operatorname{Cov}(X,Y)=\beta\) for Fig.~1(a)
and \(\operatorname{Cov}(X,Y)=\beta+\operatorname{Cov}(U_Y,U_X)\) for
Fig.~1(b), are the relations those rules give at once in this example
\cite[sec.~3.1, p.~105]{Pearl2009CausalOverview}. These are the
source's historical claims; the cited pages do not state the rules
themselves, and this chapter does not state them. Pearl adds that the
relations can also be derived from the equations, but that for large
models algebraic methods tend to obscure where the derived quantities come
from (p.~105). Here the derivation is algebraic: item~3 of
\cref{thm:p15-w02c-p15-009-linear-effect} gives
\(\operatorname{Cov}(X,Y)=\beta\operatorname{Var}(X)
+\operatorname{Cov}(U_X,U_Y)\) when \(X=U_X\). In this small case the
origin of each summand can still be read off the diagram, which is the
point of the graphical rules: with \(\operatorname{Var}(X)=1\), the
summand \(\beta\) belongs to the path \((X,Y)\) of item~2 of
\cref{thm:p15-w02c-p15-009-diagram-paths}, and the summand
\(\operatorname{Cov}(U_Y,U_X)\) to the path \((X,U_X,U_Y,Y)\) through
the double arrow; in Fig.~1(a) that path and its summand are both
absent, which leaves~(3). This path-by-path reading of~(4) is the
handbook's observation about the example.

Under certain conditions, Pearl continues, such relations let one solve
for the path coefficients in terms of observed covariances alone; his
example of a condition is \(\operatorname{Cov}(U_Y,U_X)=0\). One then
learns how large the direct causal effects are from associations observed
without experiment, with a proviso the source attaches: the inference
binds only a reader ``prepared to defend the causal assumptions'' that the
diagram encodes (p.~105). \Cref{thm:p15-w02c-p15-009-path-solve} gives the finite form
of both the conclusion and its proviso. Under the missing double arrow
and the missing arrow from \(Y\) to \(X\), the coefficient \(\beta\)
equals \(\operatorname{Cov}(X,Y)/\operatorname{Var}(X)\), a function
of the observational law of \((X,Y)\), and it then fixes every average
intervention effect of \(X\) on \(Y\). Three models with one common
observational law show that the law alone does not settle the two
assumptions: adding either link changes the intervention effect without
changing what is observed.

Pearl also stresses where a path diagram carries its causal assumptions:
in the missing links, not in the links (p.~105). An arrow indicates only
the possibility of a causal connection, whose strength remains to be
determined from data; a missing arrow claims zero influence, and a
missing double arrow claims zero covariance. In Fig.~1(a) the direct
effect \(\beta\) is identified because of the missing double arrow,
\(\operatorname{Cov}(U_Y,U_X)=0\), together with the missing arrow from
\(Y\) to \(X\); adding either link would relax the corresponding
assumption, and \(\beta\) could no longer be identified. Both assumptions
are, in his words, ``causal, not associational'', since neither can be
determined from the joint density of the observed \(X\) and \(Y\). The
association between \(U_Y\) and \(U_X\) can be uncovered only in an
experimental setting or, in more intricate models, from other causal
assumptions (p.~105). In \cref{def:p16-finite-scm} a missing arrow is a
parent set that omits a vertex. The absence of double arrows corresponds
to the independence hypothesis on the disturbances, which in this chapter
is mutual independence of the whole family
(\cref{def:p16-causal-kernels}), a stronger condition than the zero
covariances of the linear reading. The finite counterpart of ``causal, not
associational'' is item~1 of \cref{thm:p16-causal-nonidentification}: one
model there has independent disturbances and the other dependent ones, yet
the two have the same positive observational law of \((X,Y)\).

The central idea of the nonparametric extension, Pearl writes, is to use
what is invariant about structural equations while leaving their
functional form uncommitted \cite[sec.~3.2, p.~107]{Pearl2009CausalOverview}.
Accordingly \cref{def:p16-finite-scm} lets each \(f_i\) be an arbitrary
total function on its finite domain, and the source's nonparametric reading
of its three-variable chain is \cref{def:p16-causal-chain}. Finally, a
system of such functions is called \emph{structural} when it is assumed
autonomous, meaning that no function would change if the form of any other
function were altered. Pearl attributes this condition to Simon (1953)
and Koopmans (1953) (p.~107). It is the invariance on which the
equation-replacement operation of the next section relies.

\begin{hbtheorem}{thm:p16-causal-solution}{proved=annex:p16.causal-solution}{Existence, uniqueness and the observational law}
For every \(\omega\in\Omega\), the model has exactly one solution.
Evaluating the structural functions in any topological order gives that
solution. Denote its coordinates by \(V_i(\omega)\), and its tuple by
\(V(\omega)\). The \(V_i\) are finite random variables. Their joint
\emph{observational law} is the mass function
\[
 p_M(v)\defeq\mathbb P(V=v)
       =\sum_{\omega:V(\omega)=v}q(\omega),\qquad v\in S_N.
\]
For \(J\subseteq N\), its marginal law is
\[
 p_{M,J}(s)=\mathbb P(V_J=s)
           =\sum_{v:v_J=s}p_M(v),\qquad s\in S_J.
\]
All these laws are defined even when some outcomes or assignments have
probability zero.
\end{hbtheorem}

\section{Interventions replace specified equations}
\label{sec:p16-finite-interventions}

\begin{hbdefinition}{def:p16-causal-intervention}{A fixed intervention and its distribution}
Choose \(I\subseteq N\) and \(x\in S_I\). The intervention
\(\operatorname{do}(V_I=x)\) replaces the equation for every \(i\in I\)
by \(v_i=x_i\). It preserves \(\Omega,q\), all disturbances and
every equation indexed outside \(I\). Write \(M_{I=x}\) for this
model. Its graph deletes every edge whose terminal vertex is in \(I\).
For \(i\in I\), the new structural function is the constant \(x_i\),
with no endogenous parent; its old disturbance can be retained as an
unused input.

This operation stipulates that the unreplaced mechanisms and the
exogenous probability law remain invariant under the modeled manipulation.
That stipulation is part of the causal interpretation of the model; it
is the autonomy condition that the source calls the structural reading
of the equations (Simon 1953; Koopmans 1953), as recorded after
\cref{def:p16-finite-scm}. If an
actual manipulation also changes another mechanism or the population's
exogenous law, the displayed replacement alone does not model that change.

Write \(V^{I=x}(\omega)\) for its solution and
\(p_M^{I=x}(v)=\mathbb P(V^{I=x}=v)\) for its law. The traditional
intervention notation means
\[
 \mathbb P_M(V_J=s\mid\operatorname{do}(V_I=x))
 \defeq\mathbb P(V_J^{I=x}=s).
\]
Here the expression following the vertical bar names an operation on the
model. It is not a conditioning event. In particular this definition
requires no assumption that \(\mathbb P(V_I=x)>0\).

For an event \(D\subseteq S_J\), sum the displayed law over \(s\in D\).
For a real-valued function \(h:S_J\to\mathbb R\), the corresponding
intervention mean is the finite expectation
\(\mathbb E[h(V_J^{I=x})]\). A difference of such means at two fixed
interventions is a model's \emph{average intervention effect} for that
specified response function and contrast.

Other comparisons use the same intervention laws. For real \(h\) and
two values \(x',x\in S_I\), the \emph{experimental risk ratio} is
\[
 \mathbb E[h(V_J^{I=x'})]\big/\mathbb E[h(V_J^{I=x})],
\]
defined exactly when the denominator is nonzero. The \emph{intervention
variance} is
\(\operatorname{Var}(h(V_J^{I=x}))=\mathbb E[h(V_J^{I=x})^2]
-(\mathbb E[h(V_J^{I=x})])^2\), and any other functional of the law of
\(V_J^{I=x}\) may enter a comparison. Each of them is a function of the
\emph{controlled distribution}
\[
 s\longmapsto\mathbb P_M(V_J=s\mid\operatorname{do}(V_I=x))
 =\sum_{v:\,v_J=s}p_M^{I=x}(v),\qquad s\in S_J,
\]
the marginal on \(J\) of the post-intervention law; the sum runs over
the remaining coordinates, which play the role of covariates.
\index{intervention!equation replacement}\index{do operator}
\end{hbdefinition}

The source states these measures for a treatment \(X\) and response
\(Y\): the average difference
\(E(Y\mid do(x_0'))-E(Y\mid do(x_0))\), its equation~(8), the
experimental risk ratio \(E(Y\mid do(x_0'))/E(Y\mid do(x_0))\), its
equation~(9), and the variance \(\operatorname{Var}(Y\mid do(x_0))\) or
any other distributional parameter. All of them are obtained from the
controlled distribution
\(P(Y=y\mid do(x))=\sum_zP(z,y\mid do(x))\), which Pearl (2000a, 1995a)
called the ``causal effect''
\cite[sec.~3.2.1, p.~108]{Pearl2009CausalOverview}. The finite
definitions above are these, with \(h\) the response and \(J\) its
coordinates; the risk ratio requires a nonzero denominator, which the
source leaves implicit.

The source reads the post-intervention law, which it also calls the
``controlled'' or ``experimental'' distribution, in contrast with the
pre-intervention law of the unmodified model. If \(X\) is a treatment,
\(Y\) a response and \(Z\) a covariate affecting the treatment received,
then \(P(z,y\mid do(x_0))\) is the proportion of individuals who would
reach response level \(y\) and covariate level \(z\) in the hypothetical
situation where treatment \(x_0\) is given uniformly to the whole
population (p.~108). In the finite model this proportion is
\(\mathbb P(V^{I=x}_J=s)\), computed with the unchanged outcome weights
\(q\): every unit is retained and each receives the same treatment value.

The opening of the source's sec.~3.2 motivates the definition
\cite[sec.~3.2, p.~107]{Pearl2009CausalOverview}. Structural equation
modeling, Pearl writes, has been the main vehicle for effect analysis in
economics and in the behavioral and social sciences, citing Goldberger
(1972), Duncan (1975) and Bollen (1989). Its methodology, however, was
developed mostly for linear analysis, and until recently nothing
comparable had been devised for models with dichotomous variables or
nonlinear dependencies. Any such extension, he argues, must detach
``effect'' from its algebraic form as a coefficient in an equation and
redefine it as a general capacity to transmit changes among variables.
He lists Haavelmo (1943), Strotz and Wold (1960), Spirtes et al.\ (1993),
Pearl (1993a, 2000a) and Lindley (2002) as having proposed the extension
based on simulating hypothetical interventions in the model, and states
that it has yielded fresh methods for defining causal effects, and for
estimating them, in models that are nonlinear or nonparametric, the latter being models whose
equations have an unknown functional form. At the start of sec.~3.2.1
(p.~107) he adds that the invariance of autonomous equations is what
allows structural equations to serve as a basis for modeling both causal
effects and counterfactuals. The intervention laws of this section are the
first use; the finite same-unit counterfactuals of
\cref{sec:p16-finite-counterfactuals} are built on the same retained
equations and disturbances. For the linear models of the source's
Fig.~1, the effect coefficient that this extension generalizes is treated
in \cref{thm:p15-w02c-p15-009-linear-effect}.

\begin{hbtheorem}{thm:p16-causal-intervention-laws}{proved=annex:p16.causal-intervention-laws}{Laws of fixed equation replacement}
Every intervention above has a unique solution at every outcome, and
\[
 p_M^{I=x}(v)=\sum_{\omega:V^{I=x}(\omega)=v}q(\omega)
\]
is a mass function. It vanishes unless \(v_I=x\). In addition:
\begin{enumerate}
\item If \(V_I(\omega)=x\), then
      \(V^{I=x}(\omega)=V(\omega)\).
\item If \(j\notin I\cup\operatorname{de}_G(I)\), then
      \(V_j^{I=x}(\omega)=V_j(\omega)\) for every \(\omega\).
\item Intervening on the empty set leaves the model unchanged.
      For disjoint \(I,J\), performing the two replacements
      \(I=x\) and \(J=y\) in either order gives the same model
      as the single replacement on \(I\cup J\) with these values.
\end{enumerate}
The first two conclusions are pointwise, including outcomes of zero mass.
\end{hbtheorem}

Conditioning on \(V_I=x\), when this event has positive probability,
reweights the original outcomes to
\(q(\omega)\mathbf1_{\{V_I=x\}}(\omega)/\mathbb P(V_I=x)\)
and retains the original solution functions. Equation replacement retains
the original weights and evaluates the modified functions. The confounded
example below calculates the resulting difference explicitly.

\section{Independent disturbances and factorization}
\label{sec:p16-causal-factorization}

\begin{hbdefinition}{def:p16-causal-kernels}{Independent disturbances and structural kernels}
Put \(r_i(u)=\mathbb P(U_i=u)\). The disturbances are
\emph{mutually independent} if
\[
 \mathbb P(U_1=u_1,\ldots,U_n=u_n)=\prod_{i=1}^n r_i(u_i)
 \quad\text{for every }(u_i)_{i\in N}\in\prod_{i\in N}\mathcal U_i.
\]
This is independence of the entire family, not just of every pair.
Following the source, a model with an acyclic graph and mutually
independent disturbances is called \emph{Markovian}
\cite[sec.~3.2.3, p.~110]{Pearl2009CausalOverview}.

For any model, define its structural kernels by
\[
 k_i(b\mid a)\defeq
 \sum_{\substack{u\in\mathcal U_i\\f_i(a,u)=b}}r_i(u),
 \qquad a\in S_{\operatorname{pa}_G(i)},\quad b\in S_i.
\]
For each fixed \(a\), this is a mass function on \(S_i\): the sets
of disturbance values giving each output partition \(\mathcal U_i\).
The kernel is defined at every parent assignment, including assignments
which never occur observationally. Its vertical bar separates output from
input; it does not define a conditional probability on a null event.
\index{causal model!Markovian}\index{structural kernel}
\end{hbdefinition}

\begin{hbtheorem}{thm:p16-causal-factorization}{proved=annex:p16.causal-factorization}{The finite causal factorization}
In a model with mutually independent disturbances,
\[
 p_M(v)=\prod_{i\in N}k_i(v_i\mid v_{\operatorname{pa}_G(i)})
 \qquad(v\in S_N).
\]
If \(\mathbb P(V_{\operatorname{pa}_G(i)}=a)>0\), then
\[
 k_i(b\mid a)=
 \mathbb P(V_i=b\mid V_{\operatorname{pa}_G(i)}=a).
\]
Thus the structural kernels supply the observational conditional factors
on every parent event on which conditioning is defined. The product formula
itself requires no positivity assumption.
\end{hbtheorem}

\begin{hbtheorem}{thm:p16-causal-truncation}{proved=annex:p16.causal-truncation}{Truncated factorization for simultaneous interventions}
In the same model, for \(I\subseteq N\), \(x\in S_I\), and
\(v\in S_N\),
\[
 p_M^{I=x}(v)=
 \begin{cases}
 \displaystyle\prod_{i\notin I}k_i(v_i\mid v_{\operatorname{pa}_G(i)}),
       &v_I=x,\\[1ex]
 0,    &v_I\ne x.
 \end{cases}
\]
In the first case the values in intervened parent positions are therefore
the specified components of \(x\). Marginal intervention probabilities
are obtained by summing this mass function over the remaining coordinates.
For \(I=N\) the empty product is one, giving point mass at \(x\).
\end{hbtheorem}

This is a fully defined finite version of the source's causal Markov
condition and truncated factorization
\cite[Theorem~1 and Corollary~1, pp.~110--111]{Pearl2009CausalOverview}.
Using observational conditional probabilities in place of the kernels
also requires that the relevant conditioning events have positive mass.
An intervention can assign a parent configuration that was absent from
the observational distribution.

The source's note to these two results records their history
\cite[p.~111, n.~8]{Pearl2009CausalOverview}. Pearl (2000a, p.~30)
contains a short proof of the causal Markov theorem. The theorem was
first presented by Pearl and Verma (1991), and Pearl describes it as
implicit in the work of Kiiveri et al.\ (1984) and others. The truncated
factorization was named the ``Manipulation Theorem'' by Spirtes et
al.\ (1993), and Pearl describes it as implicit in Robins's (1987)
G-computation formula. The note also refers the reader to Lauritzen
(2001). These are the source's attributions; the proofs above do not
depend on them.

\begin{hbtheorem}{thm:p15-w02c-p15-009-five-variable}{proved=annex:p15-w02c-p15-009-five-variable}{The source's five-variable Markovian model}
Let the endogenous variables be \(Z_1,Z_2,Z_3,X,Y\), with the DAG whose
edges are
\[
\begin{gathered}
 Z_1\to Z_3,\quad Z_2\to Z_3,\quad Z_1\to X,\quad Z_3\to X,\\
 Z_2\to Y,\quad Z_3\to Y,\quad X\to Y,
\end{gathered}
\]
and mutually independent disturbances. Write the kernels as
\(k_1(z_1)\), \(k_2(z_2)\), \(k_3(z_3\mid z_1,z_2)\),
\(k_X(x\mid z_1,z_3)\) and \(k_Y(y\mid z_2,z_3,x)\). Then:
\begin{enumerate}
\item For every assignment,
      \[
      \begin{aligned}
       p_M(z_1,z_2,z_3,x,y)
       &=k_1(z_1)k_2(z_2)k_3(z_3\mid z_1,z_2)\\
       &\qquad\times k_X(x\mid z_1,z_3)k_Y(y\mid z_2,z_3,x).
      \end{aligned}
      \]
\item For every \(x_0\),
      \[
      \begin{aligned}
       &\mathbb P_M(Z_1=z_1,Z_2=z_2,Z_3=z_3,Y=y\mid
                    \operatorname{do}(X=x_0))\\
       &\qquad=k_1(z_1)k_2(z_2)k_3(z_3\mid z_1,z_2)
               k_Y(y\mid z_2,z_3,x_0),
      \end{aligned}
      \]
      and summing over \(z_1,z_2,z_3\) gives
      \[
      \begin{aligned}
       &\mathbb P_M(Y=y\mid\operatorname{do}(X=x_0))\\
       &\qquad=\sum_{z_1,z_2,z_3}k_1(z_1)k_2(z_2)k_3(z_3\mid z_1,z_2)
                         k_Y(y\mid z_2,z_3,x_0).
      \end{aligned}
      \]
\item For every \(x\) and \(z_2\),
      \[
      \begin{aligned}
       &\mathbb P_M(Z_1=z_1,Z_3=z_3,Y=y\mid
                    \operatorname{do}(X=x,Z_2=z_2))\\
       &\qquad=k_1(z_1)k_3(z_3\mid z_1,z_2)k_Y(y\mid z_2,z_3,x),
      \end{aligned}
      \]
      and
      \[
      \begin{aligned}
       &\mathbb P_M(Y=y\mid\operatorname{do}(X=x,Z_2=z_2))\\
       &\qquad=\sum_{z_1,z_3}k_1(z_1)k_3(z_3\mid z_1,z_2)
                k_Y(y\mid z_2,z_3,x).
      \end{aligned}
      \]
\item Let \(\mathcal Z^+\) be the set of \((z_1,z_2,z_3)\) with
      \(\mathbb P(Z_1=z_1,Z_2=z_2,Z_3=z_3)>0\), and abbreviate
      \(\mathbb P(z_3\mid z_1,z_2)=
        \mathbb P(Z_3=z_3\mid Z_1=z_1,Z_2=z_2)\) and
      \(\mathbb P(y\mid z_2,z_3,x)=
        \mathbb P(Y=y\mid Z_2=z_2,Z_3=z_3,X=x)\).
      If \(\mathbb P(Z_2=z_2,Z_3=z_3,X=x_0)>0\) for every
      \((z_1,z_2,z_3)\in\mathcal Z^+\), then
      \[
      \begin{aligned}
       &\mathbb P_M(Y=y\mid\operatorname{do}(X=x_0))\\
       &\qquad=\sum_{(z_1,z_2,z_3)\in\mathcal Z^+}\mathbb P(Z_1=z_1)
        \mathbb P(Z_2=z_2)\mathbb P(z_3\mid z_1,z_2)
        \mathbb P(y\mid z_2,z_3,x_0).
      \end{aligned}
      \]
      If \(\mathbb P(Z_2=z_2)>0\) and
      \(\mathbb P(Z_2=z_2,Z_3=z_3,X=x)>0\) for every \((z_1,z_3)\) with
      \((z_1,z_2,z_3)\in\mathcal Z^+\), then
      \[
      \begin{aligned}
       &\mathbb P_M(Y=y\mid\operatorname{do}(X=x,Z_2=z_2))\\
       &\qquad=\sum_{(z_1,z_3):\,(z_1,z_2,z_3)\in\mathcal Z^+}
        \mathbb P(Z_1=z_1)\mathbb P(z_3\mid z_1,z_2)
        \mathbb P(y\mid z_2,z_3,x).
      \end{aligned}
      \]
      Every conditional probability displayed has a conditioning event
      of positive probability.
\end{enumerate}
Without these positive-mass conditions the kernel expressions of items
2 and~3 remain exact, but an observational conditional probability at a
null conditioning event is undefined, so the observational formulas of
item~4 are not available there.
\end{hbtheorem}

Pearl presents this model as his Fig.~3, with the factorization of
item~1 as equation~(18), the post-intervention law of item~2 as~(19)
and the effect as~(20) \cite[sec.~3.2.3, pp.~111--112]{Pearl2009CausalOverview}.
He introduces the model to show what the truncated factorization is for
(p.~112): it lets causal quantities be derived directly, without handling
the equations or modifying them as in his derivation~(11)--(13). For
Fig.~3, whose error variables are left implicit, one does not write the
five nonparametric equations at all; one writes the joint distribution~(18)
directly as a product of one factor per variable, then truncates it
by~(17) to obtain~(19). In the finite formulation this is the passage from
\cref{thm:p16-causal-factorization} to \cref{thm:p16-causal-truncation}:
once the kernels are known, items~1 and~2 of
\cref{thm:p15-w02c-p15-009-five-variable} are computed from
them alone, without solving the structural equations, while the
equation-level derivation of \cref{thm:p15-w02c-p15-009-chain-weak} works
with the functions and disturbances themselves. Pearl makes two further
claims about the model. Each factor on the right of~(18) is
directly estimable from data; item~4 records the positive-mass
conditions under which an observational conditional probability equals
the corresponding kernel. And (20) is exactly what is commonly called
adjusting for \(Z_1\), \(Z_2\) and \(Z_3\), a formula that can be written
by inspection from the truncated product and marginalization, without
first deciding whether the \(Z\) variables are confounders or lie on
causal pathways; such questions can be answered from the graph, but the
truncated factorization settles them automatically.
Item~3 is the source's simultaneous or sequential intervention
\(\operatorname{do}(X=x),\operatorname{do}(Z_2=z_2)\), with
treatments \(X,Z_2\) and measured covariates \(Z_1,Z_3\): its displays
are the source's equations~(21) and~(22) (p.~112). Pearl notes that the
truncated factorization is not restricted to a single intervened
variable and applies to the analysis of time-varying treatments with
time-varying confounders, citing Robins (1986) and Arjas and Parner
(2004). In the finite formulation the simultaneous intervention is one
replacement on \(\{X,Z_2\}\), which by
\cref{thm:p16-causal-intervention-laws} equals the two replacements
performed in either order.

Two further remarks close the source's discussion
\cite[sec.~3.2.3, pp.~112--113]{Pearl2009CausalOverview}. Expression~(22)
coincides with Robins's (1987) G-computation formula, which Robins
derived from a more complicated set of counterfactual assumptions. And,
as Robins noted, the formula prescribes adjustment for covariates that may
be affected by earlier treatments, the example being \(Z_3\), which may be
affected by \(Z_2\). Item~3 shows this exactly: \(Z_2\to Z_3\) is an
edge, so \(Z_3\) is a descendant of the treatment \(Z_2\), yet the sum
in~(22) runs over \(z_3\) with weight \(k_3(z_3\mid z_1,z_2)\) evaluated
at the set treatment value \(z_2\), and \(z_3\) also enters the response
factor \(k_Y(y\mid z_2,z_3,x)\), \(Z_3\) being a parent of \(Y\). The
truncated product therefore keeps the factor of a post-treatment covariate
and deletes only the factors of the intervened vertices. This complements
the general remark of the source's note~8, recorded above, that the
truncated factorization is implicit in the G-computation formula.

\section{Identification from a specified observational law}
\label{sec:p16-finite-causal-identification}

\begin{hbdefinition}{def:p16-causal-identification}{Identification relative to a model class}
Let \(\mathcal C\) be a class of models with common observed coordinates
\(O\) and value sets \(S_O\), and let \(Q(M)\) be a specified
quantity associated with each model in the class. The quantity is
\emph{identified from the observational law on \(O\)} within
\(\mathcal C\) if, for every \(M_1,M_2\in\mathcal C\),
\[
 p_{M_1,O}=p_{M_2,O}\quad\Longrightarrow\quad Q(M_1)=Q(M_2).
\]
Equality of laws means equality at every assignment in \(S_O\).
A pair with equal observed laws and unequal values of \(Q\) proves
\emph{nonidentification} in that class. This is the two-model criterion
of the source's Definition~2 \cite[p.~109]{Pearl2009CausalOverview}.
Identification treats the observational law as given exactly; recovering
it or the target from finite sampled data is a separate estimation problem.
\index{causal identification}\index{nonidentification}
\end{hbdefinition}

The source puts identification at the centre of the analysis of causal
effects \cite[sec.~3.2.1, p.~108]{Pearl2009CausalOverview}. The question
is whether the controlled law that follows an intervention,
\(P(Y=y\mid do(x))\), can be estimated from data generated under the law
that holds before it, \(P(z,x,y)\).
\Cref{def:p16-causal-identification} makes this question exact by taking
the pre-intervention law itself as given. Within a stated class, the
controlled distribution can be recovered from data governed by the
observational law only if it is determined by that law, that is, only if
models with equal observational laws have equal controlled distributions;
and when it is so determined it is a function of the observational law,
by the equivalence proved after the definition's gloss below. The further
step of recovering the observational law, and with it the target, from
finitely many sampled observations is the separate estimation problem
mentioned in the definition. The handbook keeps the two steps apart; the
source's question, as quoted, concerns the passage from the
pre-intervention distribution to the controlled one.

The source attributes its Definition~2 to Pearl (2000a, p.~77) and

glosses it as follows \cite[sec.~3.2.1, pp.~108--109]{Pearl2009CausalOverview}.
The particular details of \(M_1\) and \(M_2\) are irrelevant. What
matters is that the assumptions defining the class, for example those
encoded in the diagram, restrict how those details can vary, so tightly
that equal observed laws force equal values of \(Q\). When this holds,
\(Q\) depends on the observed law alone and should be expressible through
its parameters. In the finite formulation this is exact: \(Q\) is
identified within \(\mathcal C\) precisely when there is a function
\(\Phi\), defined on the set of observed laws \(p_{M,O}\) with
\(M\in\mathcal C\), such that \(Q(M)=\Phi(p_{M,O})\) for every
\(M\in\mathcal C\). Given identification, put
\(\Phi(p)=Q(M)\) for any \(M\in\mathcal C\) with \(p_{M,O}=p\); the
implication in the definition makes this value independent of the
choice of \(M\). Conversely such a \(\Phi\) gives
\(Q(M_1)=\Phi(p_{M_1,O})=\Phi(p_{M_2,O})=Q(M_2)\) whenever the laws
agree.

The source motivates the definition by a contrast. In econometrics
(Hurwicz 1950; Marschak 1950; Koopmans 1953) and in social science
(Duncan 1975; Bollen 1989), identification was usually posed in linear
parametric settings, where it asks whether each parameter of the model
is determined uniquely by the parameters of the observed
distribution. A nonparametric causal quantity such as
\(Q(M)=P(y\mid do(x))\) has no distinct parametric signature: it is
defined procedurally, by simulating an intervention in the model, so
the notion of a unique solution does not apply to it directly, and the
two-model criterion replaces it.

\begin{hbtheorem}{thm:p16-causal-positive-identification}{proved=annex:p16.causal-positive-identification}{Identification with full observation and positive support}
Fix the DAG and finite value sets. Consider the class of models with
mutually independent disturbances in which all endogenous variables are
observed and \(p_M(v)>0\) for every \(v\in S_N\).
For every fixed intervention \(I=x\), its full distribution is identified
in this class. Explicitly, for \(v_I=x\), it is
\[
 p_M^{I=x}(v)=\prod_{i\notin I}
 \frac{\displaystyle\sum_{w:w_i=v_i,\,
                      w_{\operatorname{pa}_G(i)}=v_{\operatorname{pa}_G(i)}}p_M(w)}
      {\displaystyle\sum_{w:w_{\operatorname{pa}_G(i)}=v_{\operatorname{pa}_G(i)}}p_M(w)}.
\]
It is zero for \(v_I\ne x\). Every denominator is positive.
Consequently all its marginals and finite response means are identified.
Full positive support is a sufficient hypothesis here; no claim that it is
necessary for every particular target is made.
\end{hbtheorem}

The source states the corresponding claim without qualification. Causal
analysis in graphical models, Pearl writes, begins by recognizing that
every causal effect is identifiable in a Markovian model, one whose graph
is acyclic and whose error terms are all jointly independent \cite[sec.~3.2.3, p.~110]{Pearl2009CausalOverview}.
Non-Markovian models, for example those with correlated errors produced by
unmeasured confounders, allow identification only under certain
conditions, which can also be determined from the graph structure; the
source develops them in its sec.~3.3. The closing paragraph of
sec.~3.2.2 (p.~110) adds that, since graphs encode all the information
that nonparametric structural equations represent, the identification
problem should be solvable without algebraic analysis.

\Cref{thm:p16-causal-positive-identification} is the finite qualification
of that unqualified claim, and each of its two hypotheses has a role.
Full observation: the source's Theorem~1 factorizes the law of all the
endogenous variables \(V_1,\ldots,V_n\) of the model, so reading every
factor from the observational law presupposes that all of them are
observed. If an endogenous common cause is not observed, the remaining
variables are related as by correlated errors, which is the non-Markovian
situation; \cref{def:p16-causal-latent-projection} exhibits this for the
common-cause model with \(Z\) unobserved, and item~1 of
\cref{thm:p16-causal-nonidentification} shows that the effect is then not
identified even under positive support. Positive support: in the explicit
formula each factor is an observational conditional probability whose
denominator is the probability of a parent configuration, and the
hypothesis makes every such denominator positive. Without it the claim
fails even for a fully observed Markovian model. Item~2 of
\cref{thm:p16-causal-nonidentification} gives two models on the graph
\(X\to Y\), with mutually independent disturbances and both variables
observed, whose common observational law is the point mass at
\((X,Y)=(0,0)\) but whose values of
\(\mathbb P_M(Y=1\mid\operatorname{do}(X=1))\) are zero and one. The
intervention evaluates the kernel \(k_Y(\cdot\mid1)\) at a parent value of
observational probability zero, and the observational law places no
constraint on it.

\begin{hbexample}{def:p16-causal-chain}{The source's three-variable chain in the finite formulation}
Let the endogenous variables be \(Z,X,Y\), with graph
\(Z\to X\to Y\) and mutually independent disturbances. Write their
kernels as \(k_Z(z),k_X(x\mid z),k_Y(y\mid x)\).
Then the preceding theorems give
\[
\begin{aligned}
 p_M(z,x,y)&=k_Z(z)k_X(x\mid z)k_Y(y\mid x),\\
 \mathbb P_M(Z=z,Y=y\mid\operatorname{do}(X=x))
     &=k_Z(z)k_Y(y\mid x).
\end{aligned}
\]
Summing the second expression over \(z\) gives \(k_Y(y\mid x)\),
because \(\sum_z k_Z(z)=1\). If \(\mathbb P(X=x)>0\), this equals
\(\mathbb P(Y=y\mid X=x)\) by
\cref{thm:p16-causal-factorization}. This specializes the source's
equations~(5)--(7) and (16)--(17). The stated independence and graph,
together with positive mass for the conditional expression, justify the
coincidence in this example.

Summing the second expression over \(y\) instead gives
\[
 \mathbb P_M(Z=z\mid\operatorname{do}(X=x))
 =\sum_y k_Z(z)k_Y(y\mid x)=k_Z(z)=\mathbb P(Z=z),
\]
because \(\sum_yk_Y(y\mid x)=1\) and \(k_Z(z)=\mathbb P(Z=z)\) by
\cref{thm:p16-causal-factorization}, the parent set of \(Z\) being
empty. The intervention on \(X\) leaves the law of \(Z\) unchanged.
The two factors \(k_Z(z)\) and \(k_Y(y\mid x)\) of the post-intervention
law are the same factors that appear in the pre-intervention product,
the source's equation~(16); this is the source's p.~111 calculation
\cite[sec.~3.2.3, p.~111]{Pearl2009CausalOverview}.

In the source's reading of the chain's equations~(5), the absence of
\(Z\) from the arguments of \(f_Y\) is the empirical claim that variations
in \(Z\) leave \(Y\) unchanged as long as \(U_Y\) and \(X\) stay constant
\cite[sec.~3.2, p.~107]{Pearl2009CausalOverview}. In the finite chain
this is a pointwise fact about equation replacement: under the joint
intervention \(\operatorname{do}(Z=z,X=x)\) the response is
\(f_Y(x,U_Y(\omega))\) at every outcome \(\omega\), which does not depend
on \(z\), and no intervention alters \(U_Y(\omega)\).
\end{hbexample}

The chain also illustrates the source's point that, although a single
causal assumption cannot be tested in isolation, the sum total of a
model's causal assumptions often has testable implications
\cite[sec.~3.1, p.~106]{Pearl2009CausalOverview}. In Pearl's account
the chain of his Fig.~2(a) carries seven causal assumptions, each of
them the absence of an arrow, or of a double arrow, between two of the
variables; no single one is testable on its own, yet together they
imply that \(Z\) is unassociated with \(Y\) in every stratum of \(X\)
(p.~106). The source does not list the seven. One enumeration that gives
the count, which is this handbook's reading, is the following. Between a
pair of the three variables there are three possible links, an arrow in
either direction and a double arrow, so nine in all; the chain uses two,
\(Z\to X\) and \(X\to Y\). The seven missing links are the arrows
\(X\to Z\), \(Y\to X\), \(Z\to Y\) and \(Y\to Z\), and the double
arrows between \(Z\) and \(X\), between \(X\) and \(Y\), and between
\(Z\) and \(Y\), drawn in the causal diagram of
\cref{def:p15-w02c-p15-009-causal-diagram} between the corresponding
disturbance vertices. In the finite formulation the four missing arrows
are the graph \(Z\to X\to Y\), and the three missing double arrows are
read, as throughout this chapter, as mutual independence of
\(U_Z,U_X,U_Y\). The following theorem makes the three claims exact for
arbitrary finite value sets, and shows that pairwise independence of the
disturbances would not suffice for the testable implication.

\begin{hbtheorem}{thm:p15-w02c-p15-009-chain-testable}{proved=annex:p15-w02c-p15-009-chain-testable}{The chain's testable implication and its untestable parts}
Let \(S_Z,S_X,S_Y\) be finite nonempty value sets.
\begin{enumerate}
\item \textbf{The totality implies conditional independence.} In every
      finite model with endogenous \(Z,X,Y\) on these value sets, graph
      \(Z\to X\to Y\) and mutually independent disturbances, \(Z\) and
      \(Y\) are independent given \(X\): for every \(x\) with
      \(\mathbb P(X=x)>0\) and all \(z,y\),
      \[
       \mathbb P(Z=z,Y=y\mid X=x)
       =\mathbb P(Z=z\mid X=x)\,\mathbb P(Y=y\mid X=x).
      \]
\item \textbf{The implication is testable.} If \(S_Z\) and \(S_Y\) each
      have at least two members, there is a mass function on
      \(S_Z\times S_X\times S_Y\) violating the identity of item~1 at
      some \(x\) of positive mass; it is the observational law of no
      model in item~1.
\item \textbf{No single assumption is testable.} For every mass function
      \(p\) on \(S_Z\times S_X\times S_Y\) there are finite models
      \(M_{\mathrm a}\) and \(M_{\mathrm b}\) with observational law \(p\):
      \(M_{\mathrm a}\) has graph \(Z\to X\to Y\), so it carries all four
      missing arrows, with dependent disturbances allowed;
      \(M_{\mathrm b}\) has graph \(Z\to X\), \(X\to Y\), \(Z\to Y\) and
      mutually independent disturbances, so it carries all three missing
      double arrows. Each of the seven assumptions holds in
      \(M_{\mathrm a}\) or in \(M_{\mathrm b}\), so none of them, taken
      alone, excludes any observational law.
\item \textbf{Pairwise independence does not suffice.} Let
      \(S_Z=S_X=S_Y=\{0,1\}\), \(\Omega=\{0,1\}^2\) with equal masses,
      \(U_Z(a,b)=a\), \(U_X(a,b)=b\), \(U_Y(a,b)=a\oplus b\),
      \(Z=U_Z\), \(X=U_X\) and \(Y=X\oplus U_Y\), with graph
      \(Z\to X\to Y\) (the function for \(X\) ignores \(Z\)). Each pair
      of disturbances is independent, but
      \(\mathbb P(Z=1,Y=1\mid X=0)=\tfrac12\ne\tfrac14
      =\mathbb P(Z=1\mid X=0)\,\mathbb P(Y=1\mid X=0)\).
\end{enumerate}
\end{hbtheorem}

Item~1 is the instance \(V_Z\mathrel{\perp\!\!\!\perp}V_Y\mid V_X\) of
\cref{def:p16-graphical-conditional-independence}, and it also follows
from \cref{thm:p16-finite-global-markov}: the only path between \(Z\)
and \(Y\) in the graph is \(Z\to X\to Y\), whose internal vertex \(X\)
is a noncollider in the conditioning set, so \(Z\) and \(Y\) are
d-separated given \(X\). The source reads such implications off the
diagram with d-separation (p.~106); the annex proof instead derives
item~1 directly from the factorization of
\cref{thm:p16-causal-factorization}. Item~2 shows what ``testable''
means here: a law that violates the identity refutes the conjunction of
the seven assumptions, and item~3 shows that it refutes none of them
singly. Item~4 is the finite reason for reading the missing double arrows
as mutual independence; in the source's linear reading, zero covariances
are the corresponding conditions.

Pearl states d-separation at this point as his Definition~1
\cite[sec.~3.1, p.~106]{Pearl2009CausalOverview}: a set \(S\) blocks a
path when the path has an arrow-emitting node in \(S\), or a collision
node outside \(S\) none of whose descendants lies in \(S\); a set that
blocks every path between two vertices d-separates them, and the two are
then independent given \(S\). An arrow-emitting node of a path is a
noncollider on it and a collision node is a collider, so blocking in his
sense is the failure of activity in \cref{def:p16-active-path}. He
illustrates the definition on the diagram of his Fig.~2(a), in which the
disturbances \(U_Z,U_X,U_Y\) are vertices (p.~106). The path
\(U_Z\to Z\to X\to Y\) is blocked by \(\{Z\}\) and by \(\{X\}\), each
of which emits an arrow along it, and he infers two conditional
independencies that hold in every distribution the model can generate,
however the arrows are parametrized. They are printed as
\(U_X\mathrel{\perp\!\!\!\perp}Y\mid Z\) and
\(U_Z\mathrel{\perp\!\!\!\perp}Y\mid X\). The path he names starts at
\(U_Z\), and its blocking by \(\{Z\}\) gives
\(U_Z\mathrel{\perp\!\!\!\perp}Y\mid Z\); the printed first statement
fails in some models of the chain, by item~3 of the theorem below. Next,
the path \(U_Z\to Z\to X\leftarrow U_X\) is blocked by the empty set but
not by \(\{Y\}\), because \(Y\) is a descendant of the collider \(X\);
so \(U_Z\) and \(U_X\) are independent, while their independence given
\(Y\) may hold or fail. Pearl connects this treatment of colliders with
Berkson's paradox, which he attributes to Berkson (1946): observing a
common consequence of two independent causes makes the causes dependent.
His example is a pair of independent coins whose outcomes become
dependent once one is told that at least one of them shows a tail
(p.~106). Item~4 below is this example placed in the chain, and
\cref{thm:p16-collider-numerical-example} computes the same conditional
table for two independent fair bits and their common effect.

Pearl closes the illustration with a methodological claim
(p.~106). In his view the conditional independencies that d-separation
induces are the main route by which the assumptions of a structural
equation model can be confronted with nonexperimental data: almost all
statistical tests that could invalidate the model are entailed by them.
His note~7 adds that some graphs with correlated errors have further
implications, which Shpitser and Pearl (2008) call dormant independence
(p.~106, n.~7). These are the source's claims and attribution. This
chapter neither states nor proves the dormant-independence implications,
which the source ties to graphs with correlated errors, whereas the
chain's diagram has no double arrow; and the general claim about almost
all tests is not a theorem here. For the
chain itself, item~5 below proves an exact finite counterpart: the
conditional independence of item~1 of
\cref{thm:p15-w02c-p15-009-chain-testable} is the only restriction that
the Markovian chain places on the observational law of \(Z,X,Y\), so
every observational test of that model is a test of this one implication.

\begin{hbtheorem}{thm:p15-w02c-p15-009-chain-diagram}{proved=annex:p15-w02c-p15-009-chain-diagram}{d-separation in the chain's diagram, Berkson's coins and the completeness of the chain's test}
Let \(M\) be a finite model with endogenous \(Z,X,Y\), graph
\(Z\to X\to Y\) and mutually independent disturbances
\(U_Z,U_X,U_Y\). Let \(D\) be its causal diagram
(\cref{def:p15-w02c-p15-009-causal-diagram}) with
\(E^{\leftrightarrow}=\varnothing\), regarded as a DAG on the six vertices
\(U_Z,U_X,U_Y,Z,X,Y\) with the five arrows \(U_Z\to Z\), \(U_X\to X\),
\(U_Y\to Y\), \(Z\to X\), \(X\to Y\).
\begin{enumerate}
\item The six random variables \(U_Z,U_X,U_Y,Z,X,Y\) have, on the
      product of their value sets, the joint mass function
      \[
      \begin{aligned}
       &r_Z(u_Z)\,r_X(u_X)\,r_Y(u_Y)\\
       &\quad\times\mathbf 1[z=f_Z(u_Z)]\,\mathbf 1[x=f_X(z,u_X)]\,
       \mathbf 1[y=f_Y(x,u_Y)],
      \end{aligned}
      \]
      a product of kernels along \(D\) in the sense of
      \cref{thm:p16-ancestral-marginal}. Hence every d-separation in
      \(D\) yields the corresponding conditional independence of
      \cref{def:p16-graphical-conditional-independence} among these six
      variables.
\item In \(D\), \(\{U_Z\}\) and \(\{Y\}\) are d-separated given
      \(\{Z\}\) and given \(\{X\}\), and \(\{U_Z\}\) and \(\{U_X\}\)
      are d-separated given \(\varnothing\); but \(\{U_Z\}\) and
      \(\{U_X\}\) are not d-separated given \(\{Y\}\), and \(\{U_X\}\)
      and \(\{Y\}\) are not d-separated given \(\{Z\}\). Consequently,
      in every such \(M\),
      \[
       U_Z\mathrel{\perp\!\!\!\perp}Y\mid Z,\qquad
       U_Z\mathrel{\perp\!\!\!\perp}Y\mid X,\qquad
       U_Z\mathrel{\perp\!\!\!\perp}U_X .
      \]
\item \textbf{The printed first statement fails.} Let
      \(S_Z=S_X=S_Y=\{0,1\}\), \(\Omega=\{0,1\}\) with equal masses,
      \(\mathcal U_Z=\mathcal U_Y=\{0\}\), \(U_Z=U_Y=0\),
      \(U_X(\omega)=\omega\), \(f_Z=0\), \(f_X(z,u)=u\) and
      \(f_Y(x,u)=x\). Then \(\mathbb P(Z=0)=1\) and
      \(\mathbb P(U_X=1,Y=1\mid Z=0)=\tfrac12\ne\tfrac14
      =\mathbb P(U_X=1\mid Z=0)\,\mathbb P(Y=1\mid Z=0)\).
\item \textbf{Berkson's two coins.} Let \(S_Z=S_X=S_Y=\{0,1\}\),
      \(\Omega=\{0,1\}^2\) with equal masses, and code a tail as \(1\).
      Put \(U_Z(a,b)=a\), \(U_X(a,b)=b\), \(\mathcal U_Y=\{0\}\),
      \(U_Y=0\), \(f_Z(u)=u\), \(f_X(z,u)=\max(z,u)\) and
      \(f_Y(x,u)=x\), so that \(X=Y=1\) is the report that at least one
      coin shows a tail. Then \(U_Z\) and \(U_X\) are independent, but
      \[
      \begin{aligned}
       \mathbb P(U_Z=1,U_X=1\mid Y=1)&=\tfrac13,\\
       \mathbb P(U_Z=1\mid Y=1)\,\mathbb P(U_X=1\mid Y=1)&=\tfrac49.
      \end{aligned}
      \]
      If instead \(f_X(z,u)=u\), with all else unchanged, then
      \(U_Z\mathrel{\perp\!\!\!\perp}U_X\mid Y\).
\item \textbf{The chain's test is complete.} A mass function \(p\) on
      \(S_Z\times S_X\times S_Y\), for finite nonempty value sets, is the
      observational law of some finite model with graph \(Z\to X\to Y\)
      and mutually independent disturbances if and only if, for every
      \(x\) with \(p_X(x)=\sum_{z,y}p(z,x,y)>0\) and all \(z,y\),
      \[
       p(z,x,y)\,p_X(x)=p_{ZX}(z,x)\,p_{XY}(x,y),
      \]
      where \(p_{ZX}\) and \(p_{XY}\) are the two-coordinate marginals.
      This is the identity of item~1 of
      \cref{thm:p15-w02c-p15-009-chain-testable}, multiplied through by
      \(p_X(x)^2\).
\end{enumerate}
\end{hbtheorem}

Items~3 and~4 are stipulated models. Item~3 shows that the conditional
independence of \(U_X\) and \(Y\) given \(Z\), as printed in the source,
is not implied by the chain: the path \(U_X\to X\to Y\) is active given
\(\{Z\}\). Item~4 shows both halves of the source's ``may or may not
hold'': the same graph and the same independent disturbances give
dependence or independence given \(Y\) according to the structural
function for \(X\). Item~5 concerns the observational law of the
endogenous variables only; it says nothing about the disturbances, which
are not observed, and nothing about models with correlated errors.

Pearl draws two conclusions from the p.~111 calculation
\cite[sec.~3.2.3, p.~111]{Pearl2009CausalOverview}. First, the causal
assumptions of the model let the post-intervention law be predicted from
the pre-intervention law, and since \(P(y\mid x_0)\) is estimable from
nonexperimental data, so is the causal effect of \(X\) on \(Y\). Second,
the derivation makes no assumption about how the equations are shaped
or how the error terms are distributed; the source credits the graph's
structure alone, specifically ``the identity of X's parents'', with
making it go through. In the finite chain the derivation indeed uses no property of
\(f_Z,f_X,f_Y\) beyond totality and no property of the disturbance laws
beyond their mutual independence, which is part of the source's
Markovian hypothesis. The positive-mass condition
\(\mathbb P(X=x)>0\) stated in the example is what estimability from the
observational law requires: at a value \(x\) of observational probability
zero the conditional probability \(\mathbb P(Y=y\mid X=x)\) is undefined,
and the kernel \(k_Y(\cdot\mid x)\) is unconstrained by the observational
law, as item~2 of \cref{thm:p16-causal-nonidentification} shows. The
sufficient hypotheses of the example are therefore the source's
graph-only claim made exact for a finite model.

The source's derivation of \(E(Y\mid do(x_0))=E(Y\mid X=x_0)\) in this
chain, equations~(11)--(13), rests on two assumptions, as Pearl observes:
\(Y\) is a function of \(X\) and \(U_Y\) only, and \(U_Y\) is independent
of \(\{U_Z,U_X\}\), hence of \(X\)
\cite[sec.~3.2.2, p.~109]{Pearl2009CausalOverview}. Full mutual
independence of the three disturbances is not needed; \(U_Z\) and \(U_X\)
may be dependent. Pearl compares the second assumption with the
``orthogonality'' condition \(\operatorname{Cov}(X,U_Y)=0\) of linear
models, a condition that, he notes, has routinely and often
``thoughtlessly'' been invoked to justify estimating structural
coefficients by regression (p.~109). The source ties the two assumptions to
the derivation of~(13) only. Of \(P(y\mid do(x))=P(y\mid x)\) it says
just that a similar but somewhat more involved derivation establishes it
(p.~109). That this second identity also needs only the two assumptions is
this handbook's statement, proved in the following theorem together with
the finite form of~(13).

\begin{hbtheorem}{thm:p15-w02c-p15-009-chain-weak}{proved=annex:p15-w02c-p15-009-chain-weak}{The chain coincidence under the two source assumptions}
Let a finite model have endogenous \(Z,X,Y\), graph
\(Z\to X\to Y\), and disturbances \(U_Z,U_X,U_Y\), not assumed mutually
independent. Suppose only that \(U_Y\) is independent of the pair
\((U_Z,U_X)\):
\[
 \mathbb P(U_Y=u,\ (U_Z,U_X)=w)=\mathbb P(U_Y=u)\,\mathbb P((U_Z,U_X)=w)
\]
for all values \(u,w\). Then for every \(x\) with
\(\mathbb P(X=x)>0\) and every \(y\),
\[
 \mathbb P_M(Y=y\mid\operatorname{do}(X=x))=\mathbb P(Y=y\mid X=x),
\]
and, for every real \(h\) on \(S_Y\),
\(\mathbb E[h(Y^{X=x})]=\mathbb E[h(Y)\mid X=x]\). Both sides of the
first identity equal \(\mathbb P(f_Y(x,U_Y)=y)\). Moreover
\(\mathbb P_M(Z=z,Y=y\mid\operatorname{do}(X=x))
=\mathbb P(Z=z)\,\mathbb P(Y=y\mid X=x)\).
\end{hbtheorem}

Pearl then applies derivation~(11) to the linear models of his Fig.~1,
obtaining
\[
 E(Y\mid do(x_0))=E(f_Y(x_0,u_Y))=E(\beta x_0+u_Y)=\beta x_0,
\]
his equation~(14) \cite[sec.~3.2.2, p.~110]{Pearl2009CausalOverview}.
He attaches four claims to it. The equality gives \(\beta\) its causal
meaning as an effect coefficient: in a structural, as opposed to a
regressional, model \(\beta\) is not ``interpreted'' as an effect
coefficient but ``proven'' to be one by the derivation. The coefficient
keeps this meaning however \(X\) is actually selected, through the
function \(f_X\) as in Fig.~2(a), and whether \(U_X\) and \(U_Y\) are
correlated, as in Fig.~1(b), or not, as in Fig.~1(a). Correlations can
only impede the estimation of \(\beta\) from nonexperimental data; they do
not change its definition in~(14). And structural equations, in his words,
``say absolutely nothing about the conditional expectation
\(E(Y\mid X=x)\)''; such a connection may exist in special circumstances,
for instance when \(\operatorname{cov}(X,U_Y)=0\) as in~(13), but is
otherwise irrelevant to the definition or interpretation of \(\beta\) as
an effect coefficient and to the empirical claims of~(1). Pearl presents
this statement as running against what he calls ``endless confusions in
the literature'', and sends the reader to his note~15 (p.~110). That note,
attached to the instrumental-variable discussion of sec.~3.5.2, describes
\(\beta\) as the incremental causal effect of \(X\) on \(Y\), defined by
\[
 \beta\mathrel{\stackrel{\Delta}{=}}E(Y\mid do(x_0+1))-E(Y\mid do(x_0))
 =\frac{\delta}{\delta x}E(Y\mid do(x))=\frac{\delta}{\delta x}E(Y_x),
\]
and records that attempts to give \(\beta\) a statistical interpretation
have ended in frustration, citing Holland (1988), Whittaker (1990),
Wermuth (1992), Wermuth and Cox (1993) and, for difficulties persisting
into the present century, Sobel (2008)
\cite[sec.~3.5.2, p.~124, n.~15]{Pearl2009CausalOverview}. These are the
source's historical judgment and attributions; the works named are not
consulted here. In the finite formulation the derivative is not defined,
but the difference is: by item~1 of the next theorem,
\(\mathbb E[Y^{X=x_0+1}]-\mathbb E[Y^{X=x_0}]=\beta\) whenever \(x_0\) and
\(x_0+1\) both lie in \(S_X\). The finite form follows. Its item~1 is~(14) without the standardization of the source's
note~6 and with no hypothesis on \(f_X\) or on the dependence of the
disturbances. Items~2 and~3 make the special circumstances exact. The
distinction in item~4 between zero covariance and equality of conditional
and intervention means is this handbook's, not the source's.

\begin{hbtheorem}{thm:p15-w02c-p15-009-linear-effect}{proved=annex:p15-w02c-p15-009-linear-effect}{A linear response equation and its effect coefficient}
Let a finite model have endogenous variables including \(X\) and \(Y\),
with \(\operatorname{pa}_G(Y)=\{X\}\), finite value sets
\(S_X,\mathcal U_Y\subseteq\mathbb R\), a real \(\beta\), and
\[
 f_Y(x,u)=\beta x+u\qquad(x\in S_X,\ u\in\mathcal U_Y),
\]
with \(S_Y\subseteq\mathbb R\) containing every such value. Nothing is
assumed about the other structural functions or about the joint law of
the disturbances. Then:
\begin{enumerate}
\item For every \(x_0\in S_X\),
      \(\mathbb E[Y^{X=x_0}]=\beta x_0+\mathbb E[U_Y]\). Hence for
      \(x_0,x_0'\in S_X\) the average intervention effect is
      \(\mathbb E[Y^{X=x_0'}]-\mathbb E[Y^{X=x_0}]=\beta(x_0'-x_0)\).
\item If \(\mathbb P(X=x)>0\), then
      \(\mathbb E[Y\mid X=x]=\beta x+\mathbb E[U_Y\mid X=x]\). Thus
      \(\mathbb E[Y\mid X=x]=\mathbb E[Y^{X=x}]\) exactly when
      \(\mathbb E[U_Y\mid X=x]=\mathbb E[U_Y]\), which holds in
      particular when \(U_Y\) is independent of \(X\).
\item \(\operatorname{Cov}(X,Y)=\beta\operatorname{Var}(X)
      +\operatorname{Cov}(X,U_Y)\). If \(\operatorname{Var}(X)>0\), the
      least-squares slope \(\operatorname{Cov}(X,Y)/\operatorname{Var}(X)\)
      equals \(\beta\) exactly when \(\operatorname{Cov}(X,U_Y)=0\). If
      \(X\) has no parent and \(X=U_X\), then
      \(\operatorname{Cov}(X,Y)=\beta\operatorname{Var}(X)
      +\operatorname{Cov}(U_X,U_Y)\).
\item On \(\Omega=\{-1,1\}\) with equal masses, the model
      \(X=U_X\), \(Y=\beta X+U_Y\), with \(U_X(\omega)=\omega\) and
      \(U_Y(\omega)=\gamma\omega\) for a real \(\gamma\ne0\), has
      \(\mathbb E[Y^{X=x}]=\beta x\) but
      \(\mathbb E[Y\mid X=x]=(\beta+\gamma)x\) for \(x=\pm1\). On
      \(\Omega=\{-1,0,1\}\) with equal masses, the model with
      \(U_X(\omega)=\omega\) and \(U_Y(\omega)=\omega^2-\tfrac23\) has
      \(\operatorname{Cov}(X,U_Y)=0\), so its least-squares slope is
      \(\beta\), yet
      \(\mathbb E[Y\mid X=0]=-\tfrac23\ne0=\mathbb E[Y^{X=0}]\).
\end{enumerate}
Under the source's standardization, \(\mathbb E[X]=\mathbb E[Y]=0\)
gives \(\mathbb E[U_Y]=0\) and item~1 becomes~(14); with
\(\operatorname{Var}(X)=1\), item~3 becomes the source's
equation~(4), and its equation~(3) when \(\operatorname{Cov}(U_X,U_Y)=0\).
\end{hbtheorem}

\begin{hbtheorem}{thm:p15-w02c-p15-009-path-solve}{proved=annex:p15-w02c-p15-009-path-solve}{Solving for a path coefficient from observed covariances}
Call a finite model \emph{a Fig.~1(a) model with coefficient \(\beta\)}
when its endogenous variables are \(X\) and \(Y\), its graph is
\(X\to Y\), its value sets \(S_X,S_Y,\mathcal U_X,\mathcal U_Y\) are
finite subsets of \(\mathbb R\) with \(\mathcal U_X\subseteq S_X\),
\(f_X(u)=u\), \(f_Y(x,u)=\beta x+u\), and
\(\operatorname{Cov}(U_X,U_Y)=0\). Thus such a model has the missing
arrow from \(Y\) to \(X\) and, in its linear reading, the missing double
arrow.
\begin{enumerate}
\item If \(M\) is a Fig.~1(a) model with coefficient \(\beta\) and
      \(\operatorname{Var}(X)>0\), then
      \[
       \beta=\frac{\operatorname{Cov}(X,Y)}{\operatorname{Var}(X)},
      \]
      and for \(x,x'\in S_X\),
      \(\mathbb E[Y^{X=x'}]-\mathbb E[Y^{X=x}]
      =(x'-x)\operatorname{Cov}(X,Y)/\operatorname{Var}(X)\).
      Both right sides depend on \(M\) only through the observational law
      of \((X,Y)\). Hence two such models with the same observational law
      of \((X,Y)\) and \(\operatorname{Var}(X)>0\) have the same
      coefficient and the same average intervention effects at every
      pair of values in both of their value sets for \(X\). Under the
      source's standardization \(\operatorname{Var}(X)=1\), the
      displayed identity is \(\beta=\operatorname{Cov}(X,Y)\), the
      source's equation~(3) solved for \(\beta\).
\item On \(\Omega=\{-1,1\}\) with equal masses, consider the following
      three models, each with value sets in \(\{-1,0,1\}\):
      \begin{itemize}
      \item \(M_1\): graph \(X\to Y\), \(U_X(\omega)=\omega\),
            \(U_Y(\omega)=0\), \(X=U_X\), \(Y=X+U_Y\);
      \item \(M_2\): graph \(X\to Y\), \(U_X(\omega)=U_Y(\omega)=\omega\),
            \(X=U_X\), \(Y=0\cdot X+U_Y\);
      \item \(M_3\): graph \(Y\to X\), \(U_X(\omega)=0\),
            \(U_Y(\omega)=\omega\), \(Y=U_Y\), \(X=Y+U_X\).
      \end{itemize}
      The three models share one observational law of \((X,Y)\): mass
      \(\tfrac12\) at \((-1,-1)\) and at \((1,1)\), with
      \(\operatorname{Cov}(X,Y)/\operatorname{Var}(X)=1\). The model
      \(M_1\) is a Fig.~1(a) model with coefficient \(1\), and
      \(\mathbb E[Y^{X=1}]-\mathbb E[Y^{X=-1}]=2\). In \(M_2\) the
      graph and the linear form are kept but
      \(\operatorname{Cov}(U_X,U_Y)=1\), the added double arrow; \(M_3\)
      has the arrow from \(Y\) to \(X\) instead, with mutually
      independent disturbances. In both,
      \(\mathbb E[Y^{X=1}]-\mathbb E[Y^{X=-1}]=0\).
\end{enumerate}
\end{hbtheorem}

Item~1 is the source's conclusion made exact for a finite linear
model: the direct effect of \(X\) on \(Y\) is computed from the
nonexperimental association of \(X\) and \(Y\). Item~2 is its proviso.
The computation is correct only in a model class that carries the two
causal assumptions, and the observational law does not certify that
class; the same law arises from \(M_2\) and \(M_3\), in which the
computed number is not the intervention effect. Item~1 of
\cref{thm:p16-causal-nonidentification} is the binary, nonlinear
counterpart of the role of the double arrow.

\section{A confounded association and two indistinguishable observations}
\label{sec:p16-causal-confounding-example}

\begin{hbdefinition}{def:p16-confounded-binary-model}{A binary common-cause model}
For bits \(a,b\in\{0,1\}\), let \(a\oplus b\) equal zero when
\(a=b\) and one otherwise. Take
\(\Omega=\{0,1\}^3\), with coordinates \((z,a,b)\), and mass
\[
 q(z,a,b)=\tfrac12\alpha(a)\alpha(b),\qquad
 \alpha(0)=\tfrac34,\quad\alpha(1)=\tfrac14.
\]
The coordinate disturbances \(U_Z(z,a,b)=z\),
\(U_X(z,a,b)=a\), and \(U_Y(z,a,b)=b\) are mutually independent.
Specify endogenous variables with values \(\{0,1\}\) by
\[
 Z=U_Z,\qquad X=Z\oplus U_X,\qquad Y=Z\oplus U_Y.
\]
Their DAG has edges \(Z\to X\) and \(Z\to Y\).
Call this model \(M_A\). The shared parent \(Z\) is a common cause:
its equation feeds both assignments. The following calculation makes the
confounding of the observed \(X,Y\) association precise in this model.
\end{hbdefinition}

\begin{hbtheorem}{thm:p16-causal-confounded-example}{proved=annex:p16.causal-confounded-example}{Association differs from intervention in the common-cause model}
The observational marginal of \((X,Y)\) in \(M_A\) is
\[
\begin{array}{c|cc|c}
 &Y=0&Y=1&\text{row total}\\\hline
 X=0&5/16&3/16&1/2\\
 X=1&3/16&5/16&1/2\\\hline
 \text{column total}&1/2&1/2&1
\end{array}
\]
Hence
\[
 \mathbb P(Y=1\mid X=1)=\tfrac58,\qquad
 \mathbb P(Y=1\mid X=0)=\tfrac38.
\]
For either \(x\in\{0,1\}\), however,
\[
 \mathbb P_{M_A}(Y=1\mid\operatorname{do}(X=x))=\tfrac12.
\]
With response \(Y\), the difference of the two observational conditional
means is \(1/4\), whereas the average intervention effect of setting
\(X=1\) instead of \(X=0\) is zero.

If \(Z\) is also observed, every \((X,Z)\) assignment has positive
mass and this model satisfies the explicit adjustment identity
\[
 \mathbb P_{M_A}(Y=1\mid\operatorname{do}(X=x))
 =\sum_{z=0}^1\mathbb P(Y=1\mid X=x,Z=z)\mathbb P(Z=z).
\]
Its two response factors are \(1/4\) for \(z=0\) and \(3/4\)
for \(z=1\). The identity here is proved directly from this model.
\end{hbtheorem}

\begin{hbtheorem}{thm:p16-causal-nonidentification}{proved=annex:p16.causal-nonidentification}{Equal observational laws with unequal intervention laws}
The following two nonidentification statements hold.
\begin{enumerate}
\item Among finite acyclic models of binary observed variables \(X,Y\)
      with allowed parent edge \(X\to Y\) and possibly dependent
      disturbances, the full observational law does not identify
      \(\mathbb P_M(Y=1\mid\operatorname{do}(X=1))\), even when
      every observed pair has positive probability. There are two such
      models with the table above as their common law, and target values
      \(1/2\) and \(5/8\).
\item In the class with the same binary graph and mutually independent
      disturbances, removing the positive-support hypothesis also permits
      nonidentification. There are two such models with observational
      point mass at \((X,Y)=(0,0)\), but target values zero and one
      under \(\operatorname{do}(X=1)\).
\end{enumerate}
\end{hbtheorem}

\begin{hbwarning}{def:p16-causal-latent-projection}{Dependent disturbances do not inherit the product formula}
The first common-cause example can be expressed using only endogenous
\(X,Y\): on the same \(\Omega\), set
\(\widetilde U_X=(z,a)\), \(\widetilde U_Y=(z,b)\),
\(X=z\oplus a\), and \(Y=z\oplus b\). Allow the parent edge
\(X\to Y\) while the function for \(Y\) ignores that input.
The two disturbances both contain \(z\), and are dependent: their
first coordinates agree with probability one, though each is a fair bit.
Each structural output kernel is the constant mass function \((1/2,1/2)\).
Multiplying those kernels would give \(1/4\) at every pair, contradicting
the proved \(5/16,3/16\) table. The general structural model and its
interventions remain defined; the independent-disturbance hypothesis of
\cref{thm:p16-causal-factorization} is absent.
\end{hbwarning}

\begin{hbmethod}{def:p16-causal-finite-calculation}{Calculating a specified finite intervention}
List the graph, value sets, outcome weights, disturbance functions and
structural functions. Verify that every function is total on its stated
finite domain and that the graph is acyclic. Choose the intervention and
replace its equations. Evaluate once at each \(\omega\) in a
topological order, and add \(q(\omega)\) to the mass of the resulting
assignment. Group these masses for the requested marginal or response
mean. \Cref{thm:p16-causal-solution,thm:p16-causal-intervention-laws}
prove that this finite procedure terminates and returns the specified law.

With mutually independent disturbances, one may instead calculate the
kernels and apply \cref{thm:p16-causal-truncation}. To recover kernels
from observational data, separately establish the conditional identity
and its positive denominators. State the model class when claiming an
identified effect. A table fitting observed frequencies does not provide
the structural equations, their stability under intervention, the
independence of disturbances, or measurement of an omitted variable.
\end{hbmethod}

These structural foundations establish finite acyclic equation-replacement
semantics, factorization under independent disturbances, a sufficient
identification result, and explicit nonidentification examples.
The following sections develop graphical separation and back-door
adjustment in \cref{sec:p16-graphical-separation-adjustment}, then
front-door identification in \cref{sec:p16-front-door-identification}.
Each result states its graph and support assumptions and supplies its
own proof. \Cref{sec:p16-finite-counterfactuals} develops the finite
same-unit counterfactual calculus and the binary potential-outcome bridge.
General counterfactual identification and cyclic-model semantics require
their further constructions.
